% !TeX root = ../bnuthesis-example.tex

\chapter{数学符号和公式}

\section{数学符号与计量单位}

数学符号应在全文中统一，首次出现时说明含义。本例在导言区保留了 \verb|math-style = TeX| 和 \verb|math-font = newcm| 的设置；作者应结合学科惯例统一选择符号样式，不宜在正文中反复切换。

常数和函数名可用正体表示，例如圆周率 $\uppi$、虚数单位 $\symup{i}$、函数 $\sin x$ 和自然对数 $\ln x$。微分符号使用本例定义的 \cs{dif} 命令，如 $\dif y/\dif x$；有限增量使用 $\increment x$。向量可以使用粗斜体，如 $\symbf{x}$。

计量单位可用 \pkg{siunitx} 宏包处理，以统一数值、单位和空白。例如：\SI{6.4e6}{m}、\SI{9}{\micro\meter}、\si{kg.m.s^{-1}} 和 \SIrange{10}{20}{\degreeCelsius}。这些数值只用于展示排版，不表示实验测量结果。

\section{单行与多行公式}

公式另起一行居中书写，序号置于所在行最右侧并加圆括号。引用公式时使用 \cs{eqref}，例如式~\eqref{eq:example}。式中 $x$ 为实数变量，积分上下限分别为零和圆周率。
\begin{equation}
  \int_0^{\uppi} \sin x\,\dif x = 2.
  \label{eq:example}
\end{equation}

多行公式应在运算或关系符号处换行，可使用 \env{align} 对齐。对于实数参数 $a$ 和 $b$，有
\begin{align}
  (a+b)^2
    &= a^2 + 2ab + b^2, \\
  (a-b)^2
    &= a^2 - 2ab + b^2.
\end{align}
复杂分式和多行公式应检查上下行间距，避免上下标、分式或积分限重叠。

\section{数学定理与证明}

定理环境可以使用 \pkg{amsthm} 或 \pkg{ntheorem} 宏包配置，本例保留 \pkg{amsthm}。下面的定理和证明用于演示环境的使用。

\begin{theorem}
  对任意实数 $a$ 和 $b$，均有
  \begin{equation}
    a^2+b^2 \ge 2ab.
    \label{eq:square-inequality}
  \end{equation}
  等号成立当且仅当 $a=b$。
\end{theorem}
\begin{proof}
  因为 $(a-b)^2\ge 0$，展开并移项得到式~\eqref{eq:square-inequality}。平方等于零当且仅当 $a-b=0$，故等号成立当且仅当 $a=b$。
\end{proof}

模板还提供 \env{assumption}、\env{definition}、\env{proposition}、%
\env{lemma}、\env{theorem}、\env{axiom}、\env{corollary}、\env{exercise}、
\env{example}、\env{remar}、\env{problem}、\env{conjecture} 等环境。实际写作中，定理的前提、结论及证明应完整，并明确说明所引用的已有结果。
